%% kapa: exam=1 round=19 year=2008 period=1 form=A subject=appraisal_law subject_ko=부동산관계법규 seq=2 number=42
%% source: Q-Net 감정평가사 기출문제 (cst00309) artlSeq=1001962 file=extracted/1008/1001962/2008년 제19회 감정평가사 1차시험 기출문제/1교시.hwp
%% conversion: kordoc 4.24.51 -> tools/gen1.py (auto); review=proofread-sonnet; flags=none
\begin{question}{42}
국토의 계획 및 이용에 관한 법령상 개발행위의 허가에 관한 설명으로 옳지 않은 것은? (단, 도시계획사업에 의하는 경우에는 제외함)
\begin{choices}
  \item 경작을 위한 토지의 형질변경이 아닌 공유수면의 매립을 하고자 하는 자는 특별시장․광역시장․시장 또는 군수의 허가를 받아야 한다.
  \item 건축법 제2조(정의) 제1항 제2호에 따른 건축물의 건축을 하고자 하는 자는 특별시장․광역시장․시장 또는 군수의 허가를 받아야 한다.
  \item 개발행위의 허가를 받은 자가 사업기간을 단축하는 경우에는 지체없이 그 사실을 특별시장․광역시장․시장 또는 군수에게 통지하여야 한다.
  \item 개발행위의 허가 또는 변경허가를 받지 아니하거나 사위 그 밖의 부정한 방법으로 허가 또는 변경허가를 받아 개발행위를 한 자는 1년 이하의 징역 또는 1천만원 이하의 벌금에 처한다.
  \item 특별시장․광역시장․시장 또는 군수는 기반시설의 설치 또는 그에 필요한 용지의 확보․위해방지․환경오염방지․경관․조경 등을 위하여 필요하다고 인정한 경우에는 이의 이행을 담보하기 위하여 개발행위허가를 받는 자에게 이행보증금을 예치하게 할 수 있다.
\end{choices}
\end{question}
\answer{4}
